% ============================================================
% MÓDULO: PRESENCIA DIGITAL DEL PROYECTO
% ============================================================
\clearpage
\section*{MÓDULO: PRESENCIA DIGITAL DEL PROYECTO}
\subsection*{Sitio web, redes sociales y publicidad}

\begin{recuadro}
\textbf{Objetivo del módulo:} que cada grupo publique el sitio web de su proyecto final con
dominio propio, cree sus perfiles en redes sociales y en Google, y conozca cómo funcionan las
campañas de publicidad en Google Ads y Meta Ads.
\end{recuadro}

% ------------------------------------------------------------
\section*{1. EL DOMINIO DE INTERNET}

Un \textbf{dominio} es el nombre con el que se encuentra un sitio en Internet. Funciona como
una dirección fácil de recordar que el sistema \textbf{DNS} traduce a la dirección IP del
servidor donde está alojado el sitio. Ejemplo: \texttt{miproyecto.com.ar}.

\subsubsection*{1.1. Registro en NIC Argentina}
Los dominios terminados en \texttt{.ar} se registran en \textbf{NIC Argentina}
(\texttt{nic.ar}), organismo dependiente de la Secretaría Legal y Técnica de la Nación.
\begin{itemize}
  \item \textbf{Requisitos:} CUIT o CUIL y \textbf{Clave Fiscal de ARCA nivel 2 o superior}.
        Los estudiantes menores de edad pueden hacerlo con la clave de una persona adulta
        responsable o del docente/escuela.
  \item \textbf{Zonas de registro libre:} \texttt{.com.ar}, \texttt{.ar} y \texttt{.net.ar}.
        Otras zonas tienen requisitos: \texttt{.org.ar} (organizaciones sin fines de lucro),
        \texttt{.edu.ar}, \texttt{.gob.ar}, \texttt{.tur.ar}, etc.
  \item \textbf{Costo:} se paga un arancel anual de alta y de renovación. Consultar el valor
        vigente en \texttt{nic.ar}. Si no se renueva, el dominio se da de baja y otro puede
        registrarlo.
\end{itemize}

\paragraph*{Paso a paso.}
\begin{enumerate}
  \item Ingresar a \texttt{nic.ar} con CUIT/CUIL y Clave Fiscal.
  \item En el buscador, escribir el nombre deseado, elegir la zona y ``Consultar''.
  \item Si está disponible, ``Registrar'', verificar los datos y confirmar el trámite.
  \item Pagar el arancel (tarjeta, transferencia u otros medios habilitados).
\end{enumerate}

\subsubsection*{1.2. Delegación del dominio}
Registrar el dominio no hace que aparezca un sitio: hay que \textbf{delegarlo}, es decir,
indicar en NIC Argentina cuáles son los \textbf{servidores DNS} que responderán por ese nombre.
Esos servidores los da el proveedor de alojamiento (\textit{hosting}) o un servicio de DNS
(por ejemplo, \texttt{ns1.proveedor.com} y \texttt{ns2.proveedor.com}).

\begin{enumerate}
  \item En NIC Argentina, ir al dominio $\rightarrow$ ``Delegar''.
  \item Cargar los servidores DNS que indica el hosting.
  \item Esperar la propagación (desde minutos hasta 24--48 horas).
  \item En el panel DNS del proveedor se configuran los \textbf{registros}:
        \textbf{A} (apunta a la IP del servidor), \textbf{CNAME} (alias, como \texttt{www}),
        \textbf{MX} (correo) y \textbf{TXT} (verificaciones de Google, Meta, etc.).
\end{enumerate}

Otras extensiones (\texttt{.com}, \texttt{.net}, \texttt{.io}) se compran en registradores
internacionales y se pagan en dólares.

% ------------------------------------------------------------
\section*{2. EL SITIO WEB}

\subsubsection*{2.1. Opciones para armarlo}
\begin{center}
\small
\begin{tabular}{|p{0.24\linewidth}|p{0.36\linewidth}|p{0.30\linewidth}|}
\hline
\textbf{Opción} & \textbf{Características} & \textbf{Ejemplos} \\ \hline
Constructores visuales & Arrastrar y soltar, sin programar. Rápido. Planes gratuitos con
limitaciones. & Google Sites, Wix, Odoo Website \\ \hline
Gestores de contenido (CMS) & Muy flexibles, miles de plantillas y complementos. Requieren
hosting. & WordPress \\ \hline
Tiendas en línea & Orientadas a vender productos con pagos integrados. & Tiendanube,
Shopify, WooCommerce \\ \hline
Programación propia & Control total. Se aprende HTML, CSS y JavaScript. Se puede alojar
gratis como sitio estático. & GitHub Pages, Netlify, Cloudflare Pages \\ \hline
\end{tabular}
\end{center}

\subsubsection*{2.2. Contenido mínimo del sitio del proyecto}
\begin{itemize}
  \item \textbf{Inicio:} nombre, logo, frase que explique el proyecto en una línea y una imagen.
  \item \textbf{El proyecto:} problema que resuelve, cómo funciona, fotos y videos del prototipo.
  \item \textbf{Equipo:} integrantes, roles y enlaces a LinkedIn.
  \item \textbf{Contacto:} formulario, correo, WhatsApp y redes.
  \item Diseño \textbf{adaptable a celulares} (la mayoría de las visitas llegan desde el
        teléfono), carga rápida, imágenes optimizadas y certificado \textbf{HTTPS}.
\end{itemize}

\subsubsection*{2.3. Estructura básica en HTML}
\begin{verbatim}
<!DOCTYPE html>
<html lang="es">
<head>
  <meta charset="utf-8">
  <meta name="viewport" content="width=device-width, initial-scale=1">
  <title>Mi Proyecto - E.E.S.T. N° 7</title>
  <meta name="description" content="Descripción breve para Google">
</head>
<body>
  <h1>Mi Proyecto</h1>
  <p>Qué hacemos y para quién.</p>
  <a href="https://wa.me/549XXXXXXXXXX">Contactanos</a>
</body>
</html>
\end{verbatim}

\subsubsection*{2.4. SEO: aparecer en Google}
El \textbf{SEO} (\textit{Search Engine Optimization}) son las prácticas para aparecer en los
resultados \textbf{orgánicos} (no pagos) del buscador: títulos y descripciones claros en cada
página, textos con las palabras que buscaría el cliente, buena velocidad y versión móvil, y
registrar el sitio en \textbf{Google Search Console}.

% ------------------------------------------------------------
\section*{3. REDES SOCIALES}

\begin{itemize}
  \item \textbf{Elegir pocas redes y usarlas bien:} Instagram (imágenes y videos cortos),
        Facebook (comunidades y público adulto), LinkedIn (empresas y perfil técnico),
        TikTok y YouTube (video), WhatsApp Business (atención y catálogo).
  \item \textbf{Usar cuentas profesionales o de empresa}, no personales: dan estadísticas y
        permiten hacer publicidad.
  \item \textbf{Identidad coherente:} mismo nombre de usuario, logo y colores en todas.
  \item \textbf{Calendario de publicaciones:} qué, cuándo y en qué red. Combinar contenido que
        muestra el proceso (``detrás de escena''), que enseña algo y que ofrece el producto.
  \item \textbf{Medir:} alcance, interacciones, visitas al sitio y consultas recibidas.
\end{itemize}

% ------------------------------------------------------------
\section*{4. PERFIL DE EMPRESA EN GOOGLE (GOOGLE MY BUSINESS)}

El \textbf{Perfil de Empresa en Google} (antes llamado \textit{Google My Business}) es la ficha
gratuita que aparece en Google Maps y en el buscador cuando alguien busca un negocio: dirección
o zona de cobertura, horarios, teléfono, sitio web, fotos y \textbf{reseñas} de clientes.

\begin{enumerate}
  \item Ingresar con una cuenta de Google en \texttt{google.com/business}.
  \item Cargar nombre, categoría, ubicación (o zona de servicio si no hay local) y contacto.
  \item Verificar la empresa (por video, teléfono o correo, según lo que indique Google).
  \item Completar horarios, fotos, productos y servicios, y responder siempre las reseñas.
\end{enumerate}

% ------------------------------------------------------------
\section*{5. PUBLICIDAD DIGITAL}

\subsubsection*{5.1. Google Ads}
Muestra anuncios a quien \textbf{ya está buscando} algo en Google, YouTube o Maps. Se elige:
\begin{itemize}
  \item \textbf{Palabras clave} por las que aparecer (``reparación de radios de avión'').
  \item \textbf{Ubicación} geográfica y \textbf{presupuesto diario}.
  \item \textbf{Modelo de pago:} en general se paga por clic (\textbf{CPC}).
\end{itemize}

\subsubsection*{5.2. Meta Ads (Facebook, Instagram, WhatsApp)}
Muestra anuncios según el \textbf{perfil del público} (edad, zona, intereses,
comportamientos), aunque no estén buscando el producto. Se gestiona desde el
\textbf{Administrador de anuncios} de Meta y se elige un \textbf{objetivo}: reconocimiento,
tráfico al sitio, interacción, clientes potenciales, mensajes o ventas.

\subsubsection*{5.3. Métricas básicas}
\begin{center}
\small
\begin{tabular}{|p{0.18\linewidth}|p{0.72\linewidth}|}
\hline
\textbf{Métrica} & \textbf{Significado} \\ \hline
Impresiones & Cantidad de veces que se mostró el anuncio. \\ \hline
Alcance & Cantidad de personas distintas que lo vieron. \\ \hline
Clics y CTR & Clics recibidos; CTR = clics / impresiones $\times$ 100. \\ \hline
CPC & Costo por clic = inversión / clics. \\ \hline
Conversión & Acción buscada: consulta, compra, mensaje. \\ \hline
CPA & Costo por adquisición = inversión / conversiones. \\ \hline
ROAS & Retorno de la inversión publicitaria = ventas generadas / inversión. \\ \hline
\end{tabular}
\end{center}

\begin{recuadro}
\textbf{Ejemplo:} una campaña de \$30.000 obtuvo 20.000 impresiones, 400 clics y 10
consultas. CTR = 2\%; CPC = \$75; CPA = \$3.000 por consulta. Si de esas consultas salieron
ventas por \$150.000, el ROAS es 5 (se recuperó 5 veces lo invertido).
\end{recuadro}

% ------------------------------------------------------------
\section*{6. ACTIVIDADES DEL MÓDULO}

\begin{enumerate}
  \item \textbf{Nombre y dominio:} elegir el nombre del proyecto, verificar en \texttt{nic.ar}
        que el dominio esté disponible y que no exista una marca igual.
  \item \textbf{Sitio web:} armar el sitio del proyecto final con las secciones del punto 2.2,
        publicarlo y, si se registra el dominio, delegarlo.
  \item \textbf{Redes:} crear la cuenta profesional del proyecto en al menos una red y armar un
        calendario de publicaciones de un mes.
  \item \textbf{Perfil de Empresa en Google:} redactar la ficha completa (categoría,
        descripción, servicios, fotos).
  \item \textbf{Campaña simulada:} diseñar una campaña de Google Ads o Meta Ads con objetivo,
        público, palabras clave, anuncio y presupuesto, y calcular las métricas con datos
        dados por el docente.
\end{enumerate}
